% --- Άσκηση 38931 (38931.tex) ---
\Askhsh{38931}{4}
\begin{ylh}{2}
    \item 5.3. Γεωμετρική πρόοδος
\end{ylh}

Τα φράκταλ (fractal) είναι γεωμετρικά σχήματα με ένα ιδιαίτερο χαρακτηριστικό: παρουσιάζουν αυτοομοιότητα, δηλαδή κάθε μικρό τους τμήμα μοιάζει με το σύνολο. Όσο κι αν μεγεθύνουμε ένα κομμάτι του φράκταλ, θα παρατηρούμε ότι συνεχίζει να εμφανίζει την ίδια δομή. Φράκταλ εμφανίζονται τόσο στη φύση (π.χ. στο μπρόκολο ρομανέσκο), όσο και σε δημιουργήματα του ανθρώπου (όπως έργα τέχνης ή σχήματα κατασκευασμένα με υπολογιστές). Ένα από τα πιο γνωστά φράκταλ είναι το «Τρίγωνο του Sierpinski», το οποίο κατασκευάζεται με την εξής επαναλαμβανόμενη διαδικασία:

\begin{alist}
\item Ξεκινάμε με ένα ισόπλευρο τρίγωνο πλευράς 1.
\item Συνδέουμε τα μέσα των πλευρών του και αφαιρούμε το μεσαίο τρίγωνο που σχηματίζεται.
\item Στα τρία τρίγωνα που απομένουν, επαναλαμβάνουμε την ίδια διαδικασία: εντοπίζουμε τα μέσα των πλευρών και αφαιρούμε το κεντρικό τρίγωνο.
\item Επαναλαμβάνουμε αυτό το μοτίβο επ' άπειρον.

Με τον όρο «μαύρα τρίγωνα» εννοούμε εκείνα που παραμένουν μετά από κάθε αποκοπή και δεν αφαιρούνται.

\begin{center}
\begin{tikzpicture}[scale=3]
